\appendix

% Reset section format for appendix (no РОЗДІЛ prefix)
\titleformat{\section}
  {\centering\bfseries\Large}
  {\thesection.}
  {0.5em}
  {}

% Reset ToC section format for appendix
\addtocontents{toc}{\protect\renewcommand*{\protect\cftsecpresnum}{}}
\addtocontents{toc}{\protect\setlength{\protect\cftsecnumwidth}{1.5em}}

\addcontentsline{toc}{section}{ДОДАТКИ}

% Visible ДОДАТКИ heading
\begin{center}
{\bfseries\Large ДОДАТКИ}
\end{center}
\vspace{1cm}

%% Додаток А — Обов'язковий додаток (§13 Наказу МОН №40)
\section{Список публікацій Д. С. Шепілева за темою дисертації та відомості про апробацію результатів дисертації}\label{app:publications}

\printpratsi

\setlength \tabcolsep {3pt}

\begin{center}
\begin{spacing}{1}

\begin{longtable}{|p{0.5\linewidth}|p{0.45\linewidth}|}
  \caption{Відомості про апробацію результатів дисертації}
  \label{tab:approbation} \\

\hline
\hfil \textbf{Назва заходу} &

\hfil \begin{tabular}{c}
     \textbf{Місце та дата}  \\
     \textbf{проведення}
\end{tabular}

\\

\hline

\endfirsthead

\multicolumn{2}{r}%
{{\itshape Продовження таблиці \thetable{} }} \\
\hline

\hfil \textbf{Назва заходу} &

\hfil \begin{tabular}{c}
     \textbf{Місце та дата}  \\
     \textbf{проведення}
\end{tabular}

\\

\hline

\endhead

\endfoot

\endlastfoot

3rd International Workshop on Augmented Reality in Education (AREdu 2020) & Кривий Ріг, Криворізький національний університет, 13 травня 2020~р. \\\hline

XII International Conference on Mathematics, Science and Technology Education (ICon-MaSTEd 2020) & Кривий Ріг, Криворізький державний педагогічний університет, 15--17 жовтня 2020~р. \\\hline

International Workshop on Computer Science \& Software Engineering at Schools and in the Workplace (CS\&SE@SW 2020) & Кривий Ріг, Криворізький національний університет, 27 листопада 2020~р. \\\hline

9th Workshop on Cloud Technologies in Education (CTE 2021) & Кривий Ріг, Криворізький національний університет, 17 грудня 2021~р. \\\hline

Digital Humanities Workshop (DHW 2021) & Київ, Київський університет імені Бориса Грінченка, 23 грудня 2021~р. \\\hline

11th Illia O. Teplytskyi Workshop on Computer Simulation in Education (CoSinE 2024) & Кривий Ріг, 15 травня 2024~р. \\\hline

\end{longtable}

\end{spacing}
\end{center}

%% Додаток Б
\section{Навчальний посібник з WebAR}\label{app:textbook}

Навчальний посібник <<Розробка WebAR-застосунків для освіти>> структуровано у 10 розділів, що послідовно охоплюють всі аспекти створення освітніх застосунків з доповненою реальністю у браузері:

\begin{enumerate}
    \item Вступ до WebAR
    \item Інструменти розробки (JavaScript, локальні сервери, ngrok, віддалене налагодження)
    \item Як AR працює у браузерах (WebGL, Three.js, накладання камери)
    \item Відстеження зображень (MindAR NFT, маркери, 3D-моделі, GLTF/GLB)
    \item Відстеження обличчя (орієнтири, окклюдери, маски, перемикання камери)
    \item Відстеження простору (WebXR Device API, hit testing, розміщення на поверхнях)
    \item Інтеграція машинного навчання (TensorFlow.js, handpose, face-api.js)
    \item Інші технології WebAR (A-Frame, model-viewer, 8th Wall, Zappar)
    \item Висновки
    \item Лабораторний практикум (10 варіантів творчих завдань)
\end{enumerate}

Повний текст посібника подано в електронному додатку.

%% Додаток В
\section{Лабораторний практикум}\label{app:labs}

Лабораторний практикум складається з 10 варіантів творчих завдань з 3D-моделювання та доповненої реальності. Кожний варіант включає:
\begin{itemize}
    \item завдання на створення 3D-моделей у форматі GLTF/GLB;
    \item завдання на розробку WebAR-застосунку з відстеженням зображень;
    \item завдання на розробку WebAR-застосунку з відстеженням обличчя;
    \item творче завдання на інтеграцію моделей машинного навчання у WebAR.
\end{itemize}

%% Додаток Г
\section{Структура курсу <<Синтетична реальність в освіті>> у LMS Moodle}\label{app:moodle}

Курс розгорнуто на платформі moodle.kdpu.edu.ua і структуровано у 18 тижневих модулів (лютий~-- червень 2025):

\begin{longtable}{|p{0.08\textwidth}|p{0.85\textwidth}|}
\hline
Тиждень 1 & Вступ. Реєстрація, знайомство з курсом \\
\hline
Тиждень 2 & Вступ до WebAR \\
\hline
Тиждень 3 & Інструменти розробки WebAR \\
\hline
Тиждень 4 & Як AR працює у браузерах \\
\hline
Тиждень 5 & Three.js: основи 3D-графіки \\
\hline
Тиждень 6 & Відстеження зображень (Image Tracking) \\
\hline
Тиждень 7 & Відстеження зображень: практикум \\
\hline
Тиждень 8 & Відстеження обличчя (Face Tracking) \\
\hline
Тиждень 9 & Відстеження обличчя: практикум \\
\hline
Тиждень 10 & Відстеження простору (World Tracking) \\
\hline
Тиждень 11 & WebXR Device API \\
\hline
Тиждень 12 & Інтеграція машинного навчання \\
\hline
Тиждень 13 & TensorFlow.js та Teachable Machine \\
\hline
Тиждень 14 & Інші технології WebAR \\
\hline
Тиждень 15 & A-Frame та model-viewer \\
\hline
Тиждень 16 & Віртуальна хімічна лабораторія \\
\hline
Тиждень 17 & Підсумковий проєкт \\
\hline
Тиждень 18 & Захист проєктів та рефлексія \\
\hline
\end{longtable}

%% Додаток Д
\section{Віртуальна хімічна лабораторія}\label{app:anchem}

Віртуальну хімічну лабораторію з доповненою реальністю (anchem-ar) розроблено як практичний приклад освітнього WebAR-застосунку для якісного аналізу хімічних речовин. Застосунок демонструє інтеграцію технологій MindAR та Three.js для створення інтерактивного навчального середовища.

\label{last_page}
